\documentclass[12pt]{article}
\usepackage[T1]{fontenc}
\usepackage[utf8]{inputenc}
\usepackage{lmodern}
\usepackage{textcomp}
\usepackage{enumitem}
\usepackage{geometry}
\usepackage{graphicx}
\usepackage{amsmath, amssymb}
\geometry{margin=1in}
\begin{document}
\begin{center}
Machine Learning - Undergraduate Level Assessment 18 \
Total Marks: 40 \
Answer all questions.
\end{center}

\section*{Multiple Choice (10 marks)}
\begin{enumerate}[label=\arabic*.]
\item Compared to the variance of the Maximum Likelihood Estimate (MLE), the variance of the Maximum A Posteriori (MAP) estimate is \_\_\_\_\_\_\_\_
\begin{enumerate}[label=(\alph*)]
    \item higher
    \item same
    \item lower
    \item it could be any of the above
\end{enumerate}
\item Which of the following is the joint probability of H, U, P, and W described by the given Bayesian Network H -\textgreater{} U \textless{}- P \textless{}- W? [note: as the product of the conditional probabilities]
\begin{enumerate}[label=(\alph*)]
    \item P(H, U, P, W) = P(H) * P(W) * P(P) * P(U)
    \item P(H, U, P, W) = P(H) * P(W) * P(P | W) * P(W | H, P)
    \item P(H, U, P, W) = P(H) * P(W) * P(P | W) * P(U | H, P)
    \item None of the above
\end{enumerate}
\item Predicting the amount of rainfall in a region based on various cues is a \_\_\_\_\_\_ problem.
\begin{enumerate}[label=(\alph*)]
    \item Supervised learning
    \item Unsupervised learning
    \item Clustering
    \item None of the above
\end{enumerate}
\item Statement 1| RELUs are not monotonic, but sigmoids are monotonic. Statement 2| Neural networks trained with gradient descent with high probability converge to the global optimum.
\begin{enumerate}[label=(\alph*)]
    \item True, True
    \item False, False
    \item True, False
    \item False, True
\end{enumerate}
\item Consider the Bayesian network given below. How many independent parameters would we need if we made no assumptions about independence or conditional independence H -\textgreater{} U \textless{}- P \textless{}- W?
\begin{enumerate}[label=(\alph*)]
    \item 3
    \item 4
    \item 7
    \item 15
\end{enumerate}
\end{enumerate}

\section*{True / False (10 marks)}
\textit{Answer True or False}
\begin{enumerate}[label=\arabic*.]
\setcounter{enumi}{5}
\item True or False: The rank of the matrix A = [[1, 1, 1], [1, 1, 1], [1, 1, 1]] is 1.
\item True or False: Overfitting refers to a model that cannot generalize to new data but fits the training data perfectly.
\item True or False: P(A, B, C) can be expressed as P(A, B | C) * P(C) for Boolean random variables A, B, and C.
\item True or False: Neural networks can utilize a mix of different activation functions to improve learning and performance.
\item True or False: Lasso regression is more suitable for feature selection due to its ability to shrink some coefficients to zero.
\end{enumerate}

\section*{Long Form Response (20 marks)}
\textit{Answer each question with a clear and complete explanation.}
\begin{enumerate}[label=\arabic*.]
\setcounter{enumi}{10}
\item Write a function to generate a two-dimensional array.
\item Write a function to find frequency count of list of lists.
\end{enumerate}

\vfill
\noindent\textit{End of Paper}
\end{document}